%==============================================================================
%  DevPilot AI — Rapport de Projet de Fin d'Études (ENSIAS)
%  Plateforme agentique de RAG de qualité production.
%
%  Compilation : pdflatex main.tex  (x2 pour TOC/refs).  Pas de shell-escape.
%  Les diagrammes UML sont fournis en code PlantUML (listings) + emplacement
%  d'image à rendre séparément :  plantuml diagrams/*.puml
%  Les captures d'écran sont des emplacements [CAPTURE D'ÉCRAN] à insérer.
%==============================================================================
\documentclass[12pt,a4paper,oneside]{report}

%------------------------------------------------------------------ Langue / encodage
\usepackage[utf8]{inputenc}
\usepackage[T1]{fontenc}
\usepackage[french]{babel}
\usepackage{lmodern}
\usepackage{microtype}
\usepackage{csquotes}

%------------------------------------------------------------------ Mise en page
\usepackage[a4paper,margin=2.5cm,headheight=15pt]{geometry}
\usepackage{setspace}
\onehalfspacing
\usepackage{parskip}
\usepackage{fancyhdr}
\usepackage{float}

%------------------------------------------------------------------ Tableaux / figures
\usepackage{graphicx}
\usepackage{booktabs}
\usepackage{array}
\usepackage{tabularx}
\usepackage{longtable}
\usepackage{multirow}
\usepackage{caption}
\usepackage{enumitem}
\usepackage{amsmath}
\usepackage{amssymb}

%------------------------------------------------------------------ Couleurs
\usepackage{xcolor}
\definecolor{brand}{HTML}{6A35F0}
\definecolor{brandDark}{HTML}{4C1FB3}
\definecolor{slate}{HTML}{475569}
\definecolor{codeBg}{HTML}{F6F8FA}
\definecolor{codeKw}{HTML}{8B2FB3}
\definecolor{codeStr}{HTML}{0A7E3D}
\definecolor{codeCom}{HTML}{6B7280}
\definecolor{lineGray}{HTML}{C7CDD6}

%------------------------------------------------------------------ Listings
\usepackage{listings}
\lstdefinestyle{code}{
  basicstyle=\ttfamily\footnotesize,
  keywordstyle=\color{codeKw}\bfseries,
  stringstyle=\color{codeStr},
  commentstyle=\color{codeCom}\itshape,
  numberstyle=\tiny\color{codeCom},
  backgroundcolor=\color{codeBg},
  frame=single, rulecolor=\color{lineGray},
  numbers=left, numbersep=8pt, tabsize=2,
  breaklines=true, breakatwhitespace=true, showstringspaces=false,
  columns=fullflexible, keepspaces=true, captionpos=b,
  xleftmargin=14pt, framexleftmargin=14pt,
}
\lstdefinestyle{plantuml}{
  basicstyle=\ttfamily\scriptsize,
  backgroundcolor=\color{codeBg},
  frame=single, rulecolor=\color{lineGray},
  breaklines=true, columns=fullflexible, keepspaces=true,
  showstringspaces=false, captionpos=b, xleftmargin=10pt,
}
\lstdefinestyle{eraser}{
  basicstyle=\ttfamily\scriptsize,
  backgroundcolor=\color{codeBg},
  frame=single, rulecolor=\color{lineGray},
  breaklines=true, columns=fullflexible, keepspaces=true,
  showstringspaces=false, captionpos=b, xleftmargin=10pt,
}
\lstset{style=code}

%------------------------------------------------------------------ Hyperref
\usepackage[hidelinks]{hyperref}
\hypersetup{colorlinks=true, linkcolor=brandDark, citecolor=brandDark, urlcolor=brandDark,
  pdftitle={DevPilot AI - Rapport PFE ENSIAS}, pdfauthor={DevPilot AI}}

%------------------------------------------------------------------ En-têtes
\pagestyle{fancy}
\fancyhf{}
\fancyhead[L]{\small\nouppercase{\leftmark}}
\fancyhead[R]{\small\textbf{DevPilot AI}}
\fancyfoot[C]{\thepage}
\renewcommand{\headrulewidth}{0.4pt}

%------------------------------------------------------------------ Macros
% Emplacement d'un diagramme PlantUML (le code figure dans le listing au-dessus).
\newcommand{\plantumlfig}[2]{%
  \begin{figure}[H]\centering
  \fbox{\parbox[c][3.6cm][c]{0.85\textwidth}{\centering\itshape\small\color{slate}%
  [DIAGRAMME PLANTUML]\\[3pt]#1\\[3pt]\footnotesize Rendre le code PlantUML ci-dessus (\texttt{plantuml}).}}%
  \caption{#1}\label{#2}\end{figure}}
% Emplacement d'un diagramme Eraser (le code figure dans le listing au-dessus).
\newcommand{\eraserfig}[2]{%
  \begin{figure}[H]\centering
  \fbox{\parbox[c][3.6cm][c]{0.85\textwidth}{\centering\itshape\small\color{slate}%
  [DIAGRAMME ERASER]\\[3pt]#1\\[3pt]\footnotesize Rendre le code Eraser ci-dessus sur \texttt{app.eraser.io}.}}%
  \caption{#1}\label{#2}\end{figure}}
% Emplacement d'une capture d'écran.
\newcommand{\screenshotfig}[2]{%
  \begin{figure}[H]\centering
  \fbox{\parbox[c][3.6cm][c]{0.82\textwidth}{\centering\itshape\small\color{slate}%
  [CAPTURE D'ÉCRAN]\\[3pt]#1}}%
  \caption{#1}\label{#2}\end{figure}}
\newcommand{\devpilot}{\textsc{DevPilot AI}}

%==============================================================================
\begin{document}
%==============================================================================

%------------------------------------------------------------------ Page de garde
\begin{titlepage}
\centering\thispagestyle{empty}
{\small\scshape Royaume du Maroc\\ Université Mohammed V de Rabat\\
École Nationale Supérieure d'Informatique et d'Analyse des Systèmes (ENSIAS)\par}
\vspace{0.8cm}{\color{lineGray}\rule{\textwidth}{0.5pt}}\vspace{1.0cm}

{\small\itshape Rapport de Projet de Fin d'Études\par}
\vspace{0.3cm}{\small en vue de l'obtention du diplôme d'Ingénieur d'État\par}
\vspace{0.2cm}{\small\textit{Filière : Génie Logiciel / Intelligence Artificielle}\par}

\vspace{1.4cm}
{\color{brand}\Huge\bfseries DevPilot AI\par}
\vspace{0.5cm}
{\LARGE\bfseries Conception et réalisation d'une plateforme\\[4pt]
agentique de RAG avec observabilité,\\[4pt]
traçabilité et visualisation d'exécution\par}

\vspace{1.6cm}{\color{lineGray}\rule{0.6\textwidth}{0.5pt}}\vspace{0.5cm}
\begin{minipage}{0.46\textwidth}\begin{flushleft}\large
\textbf{Réalisé par :}\\[2pt][Nom et prénom de l'étudiant]\end{flushleft}\end{minipage}\hfill
\begin{minipage}{0.46\textwidth}\begin{flushright}\large
\textbf{Encadré par :}\\[2pt][Encadrant ENSIAS]\\{[Encadrant professionnel]}\end{flushright}\end{minipage}

\vspace{1.2cm}\textbf{Membres du jury :}\\[4pt]
\begin{tabular}{ll}
[Président] & \textit{Président}\\ [Examinateur] & \textit{Examinateur}\\ [Encadrant] & \textit{Encadrant}\\
\end{tabular}

\vfill{\color{lineGray}\rule{\textwidth}{0.5pt}}\vspace{0.3cm}
{\small Année universitaire 2025--2026\par}
\end{titlepage}

%------------------------------------------------------------------ Front matter
%====================== FRONT MATTER ======================
\chapter*{Remerciements}
\addcontentsline{toc}{chapter}{Remerciements}
\markboth{Remerciements}{Remerciements}

Au terme de ce projet de fin d'études, je tiens à exprimer ma sincère gratitude
à mon encadrant pédagogique [\textit{à compléter}] pour son suivi rigoureux et
ses conseils déterminants, ainsi qu'à mon encadrant professionnel
[\textit{à compléter}] pour son expertise et sa disponibilité. Je remercie
l'ensemble du corps professoral de l'ENSIAS pour la qualité de la formation
d'ingénieur, et les membres du jury qui ont accepté d'évaluer ce travail. Je
dédie enfin ce travail à ma famille, dont le soutien a rendu ce parcours possible.

%====================== RÉSUMÉ ======================
\chapter*{Résumé}
\addcontentsline{toc}{chapter}{Résumé}
\markboth{Résumé}{Résumé}

Les équipes d'ingénierie logicielle accumulent une connaissance technique
volumineuse et fragmentée — spécifications, documents d'architecture, comptes
rendus et directives internes. La recherche par mots-clés en capture mal
l'intention sémantique, tandis que les assistants génératifs généralistes
produisent des réponses fluides mais \emph{non fondées} : hallucinations, absence
de citations et opacité du raisonnement.

\devpilot{} est une plateforme agentique de \emph{Génération Augmentée par
Récupération} (RAG) de qualité production qui répond à ces limites. Chaque
réponse est ancrée dans les documents propres à l'organisation, étayée par des
citations, évaluée automatiquement (fidélité, pertinence, risque d'hallucination)
et entièrement traçable. La génération repose sur un flux agentique de sept
agents coopérants, chacun persisté comme une exécution inspectable.

Au-delà du moteur RAG, le projet apporte trois contributions distinctives : une
tour de contrôle d'exploitation (\textbf{RAGOps}) pour l'observabilité du système
d'IA, un sous-système de \textbf{traçabilité} reconstruisant le cycle de vie
complet de chaque requête, et un \textbf{AI Execution Studio} qui visualise le
pipeline en temps réel. La plateforme est bâtie sur FastAPI, PostgreSQL, Redis,
Celery, Qdrant, Ollama et Gemini côté serveur, et Next.js côté client, le tout
orchestré par Docker Compose. Ce mémoire détaille la démarche d'ingénierie, la
conception, l'architecture et la validation du système.

\noindent\textbf{Mots-clés :} RAG agentique, grands modèles de langage, recherche
vectorielle, observabilité de l'IA, traçabilité, évaluation, FastAPI, Qdrant,
Next.js.

%====================== ABSTRACT ======================
\chapter*{Abstract}
\addcontentsline{toc}{chapter}{Abstract}
\markboth{Abstract}{Abstract}

Software engineering teams accumulate large, fragmented technical knowledge.
Keyword search misses semantic intent, while general-purpose generative
assistants produce fluent but \emph{ungrounded} answers — hallucinations, no
citations, no visibility into the reasoning.

\devpilot{} is a production-grade \emph{agentic Retrieval-Augmented Generation}
(RAG) platform that closes this gap. Every answer is grounded in the
organization's own documents, supported by citations, automatically evaluated
(faithfulness, relevance, hallucination risk) and fully traceable. Generation is
an agentic workflow of seven cooperating agents, each persisted as an inspectable
run.

Beyond the RAG engine, the project contributes three distinctive elements: an
operational control tower (\textbf{RAGOps}) for AI-system observability, a
\textbf{traceability} subsystem that reconstructs each query's full life cycle,
and an \textbf{AI Execution Studio} that visualizes the pipeline in real time.
The platform is built on FastAPI, PostgreSQL, Redis, Celery, Qdrant, Ollama and
Gemini on the server, with Next.js on the client, orchestrated by Docker Compose.
This report details the engineering process, design, architecture and validation.

\noindent\textbf{Keywords:} Agentic RAG, Large Language Models, Vector Search, AI
Observability, Traceability, Evaluation, FastAPI, Qdrant, Next.js.

%====================== INTRODUCTION GÉNÉRALE ======================
\chapter*{Introduction générale}
\addcontentsline{toc}{chapter}{Introduction générale}
\markboth{Introduction générale}{Introduction générale}

\section*{Contexte}
L'intelligence artificielle générative a transformé l'accès à l'information, mais
son adoption en entreprise se heurte à une exigence non satisfaite par les
modèles généralistes : répondre à partir du patrimoine documentaire
\emph{propre} à l'organisation, avec des garanties de véracité et de
transparence. Pour une équipe d'ingénierie, une réponse plausible mais fausse
constitue un risque opérationnel, non une simple imperfection.

\section*{Motivation}
La promesse du RAG est de réconcilier la puissance générative des grands modèles
de langage avec la connaissance vérifiable de l'organisation~\cite{lewis2020}.
Transformer cette promesse en \emph{système de production} suppose toutefois de
résoudre des problèmes d'ingénierie concrets : ingestion asynchrone, recherche
vectorielle, ancrage et citation des sources, mesure de la qualité et, surtout,
observabilité et traçabilité de bout en bout.

\section*{Problématique}
Comment concevoir et réaliser une plateforme d'IA capable de centraliser les
connaissances d'une équipe, d'en récupérer l'information pertinente, de générer
des réponses fondées et citées, d'en mesurer la qualité, et d'offrir une
traçabilité et une observabilité permettant d'auditer et d'exploiter le système
comme tout service critique d'entreprise~?

\section*{Objectifs}
Le projet vise à concevoir et implémenter \devpilot{}, plateforme RAG agentique
multi-espaces offrant : espaces de travail isolés, ingestion documentaire
asynchrone, récupération hybride avec reclassement, génération agentique ancrée
et citée, évaluation automatique, \emph{streaming} temps réel, et un sous-système
d'observabilité (RAGOps, explorateur et inspecteur de traces, AI Execution
Studio).

\section*{Méthodologie}
La démarche est itérative et orientée production : elle s'est déroulée en
\textbf{vingt-trois phases} incrémentales, de la fondation conteneurisée jusqu'aux
fonctionnalités phares de visualisation. Chaque phase fixe un objectif vérifiable,
introduit des fonctionnalités, justifie ses choix techniques et est validée par
des tests. Le chapitre~\ref{chap:archi} (architecture et réalisation) constitue
le c{\oe}ur du travail.

\section*{Organisation du rapport}
Le chapitre~\ref{chap:contexte} pose le contexte, l'analyse de l'existant et le
cahier des charges. Le chapitre~\ref{chap:conception} présente l'analyse et la
conception (UML, modèle de données). Le chapitre~\ref{chap:archi} détaille
l'architecture et la réalisation. Le chapitre~\ref{chap:techno} justifie les
technologies et l'implémentation. Le chapitre~\ref{chap:validation} traite de la
validation et des résultats. Le rapport se clôt par une conclusion générale et
des perspectives.


\tableofcontents
\listoffigures
\listoftables

%------------------------------------------------------------------ Chapitres
\chapter{Contexte général du projet}
\label{chap:contexte}

\section{Présentation générale}
Ce chapitre établit les fondations du projet \devpilot{}. Il caractérise le
besoin métier, analyse l'existant à travers une étude comparative des solutions
concurrentes, puis formalise les exigences fonctionnelles et non fonctionnelles,
les contraintes, le cahier des charges et la gestion des risques. Cette analyse
conditionne l'ensemble des décisions de conception et d'architecture présentées
dans les chapitres~\ref{chap:conception} et~\ref{chap:archi}.

\subsection{Contexte}
Les équipes d'ingénierie logicielle produisent et consomment en permanence une
connaissance technique abondante : spécifications, documents d'architecture, notes
de conception, comptes rendus, directives internes et documentation de projet.
Cette connaissance possède trois propriétés qui la rendent difficile à exploiter.
Elle est d'abord \textbf{volumineuse et croissante}. Elle est ensuite
\textbf{fragmentée}, répartie entre des outils, des formats et des silos
hétérogènes. Elle est enfin \textbf{contextuelle} : une même question n'appelle
pas la même réponse selon le projet, l'équipe ou la version concernée.

La recherche par mots-clés des outils traditionnels échoue dès que le vocabulaire
de la question diffère de celui du document ; elle ne comprend pas l'intention
sémantique et renvoie des listes de documents que l'utilisateur doit encore lire
et synthétiser. À l'inverse, les assistants génératifs généralistes produisent
une synthèse directe, mais répondent à partir de paramètres figés au moment de
l'entraînement, sans connaissance du patrimoine documentaire propre à
l'organisation et sans garantie de véracité.

\subsection{Organisme d'accueil}
Ce projet de fin d'études a été réalisé au sein de [\textit{nom de l'organisme
d'accueil — à compléter}], [\textit{secteur d'activité — à compléter}]. La
structure d'accueil [\textit{présentation succincte de l'organisme, de ses
activités et de l'équipe d'encadrement — à compléter}]. Le projet s'inscrit dans
[\textit{contexte de la mission — à compléter}].

\subsection{Problématique}
La problématique centrale se décompose en deux familles de limites
complémentaires. D'une part, les moteurs de recherche internes, fondés sur
l'appariement lexical, ne comprennent pas la sémantique des requêtes. D'autre
part, les grands modèles de langage généralistes souffrent
d'\textbf{hallucinations}, d'une \textbf{absence de traçabilité}, d'une
\textbf{absence de citations}, d'une \textbf{absence de visibilité} sur la qualité
de récupération et d'une \textbf{absence d'observabilité}~\cite{lewis2020}.

La question directrice du projet est donc la suivante : \emph{comment concevoir et
réaliser une plateforme d'IA capable de centraliser les connaissances d'une équipe
d'ingénierie, d'en récupérer l'information pertinente, de générer des réponses
fondées et citées, d'en mesurer la qualité, et d'offrir une traçabilité et une
observabilité permettant d'auditer et d'exploiter le système comme tout service
critique d'entreprise~?}

\subsection{Objectifs}
Le projet vise à concevoir et implémenter \devpilot{}, une plateforme agentique de
\emph{Génération Augmentée par Récupération} (RAG) multi-espaces. Les objectifs
opérationnels sont synthétisés dans le tableau~\ref{tab:objectifs}.

\begin{table}[H]
\centering\caption{Objectifs opérationnels et capacités correspondantes}
\label{tab:objectifs}\small
\begin{tabularx}{\textwidth}{@{}>{\bfseries}l X@{}}
\toprule
Objectif & Capacité réalisée \\
\midrule
Centraliser & Espaces de travail isolés, gestion documentaire, ingestion asynchrone \\
Récupérer & Recherche sémantique et hybride, reclassement par pertinence \\
Générer & Réponses ancrées produites par un pipeline agentique de sept agents \\
Citer & Génération de citations reliées aux passages sources \\
Mesurer & Évaluation de la fidélité, de la pertinence et du risque d'hallucination \\
Diffuser & \emph{Streaming} temps réel des réponses (Server-Sent Events) \\
Superviser & RAGOps : santé des pipelines, couverture, file de traitement, alertes \\
Tracer & Explorateur et inspecteur de traces de bout en bout \\
Visualiser & AI Execution Studio : visualisation animée du pipeline d'exécution \\
\bottomrule
\end{tabularx}
\end{table}

\subsection{Méthodologie}
La démarche adoptée est itérative et orientée production. Le développement s'est
déroulé en \textbf{vingt-trois phases} incrémentales, de la fondation
conteneurisée jusqu'aux fonctionnalités phares de visualisation. Chaque phase fixe
un objectif vérifiable, introduit des fonctionnalités, justifie ses choix
techniques et est validée par des tests dédiés. Cette progression, détaillée tout
au long du rapport, garantit que le système reflète une évolution réelle et
maîtrisée plutôt qu'une conception monolithique. Le chapitre~\ref{chap:archi}
constitue le c{\oe}ur du travail d'ingénierie.

\subsection{Analyse de l'existant}
Aucune solution unique du marché ne couvre simultanément l'ancrage documentaire
privé, la traçabilité et l'observabilité dans un cadre auto-hébergeable. Nous
analysons ci-dessous six familles de solutions représentatives.

\paragraph{ChatGPT (et assistants génériques).} \textbf{Forces :} qualité
conversationnelle, fluidité, large couverture générale. \textbf{Limites :} répond
depuis ses paramètres, pas depuis la base de connaissances privée ; hallucinations ;
absence de citations vérifiables. \textbf{Manques :} traçabilité, isolation des
données, observabilité opérationnelle.

\paragraph{LangSmith.} \textbf{Forces :} excellente observabilité et traçabilité
des applications à base de LLM~\cite{langsmith}. \textbf{Limites :} c'est un outil
d'\emph{observabilité}, non une plateforme RAG complète et auto-suffisante.
\textbf{Manques :} moteur d'ingestion/récupération intégré, interface utilisateur
métier, gouvernance multi-espaces.

\paragraph{LlamaIndex.} \textbf{Forces :} \emph{framework} riche pour
l'indexation et la récupération~\cite{llamaindex}. \textbf{Limites :} bibliothèque
de développement, non un produit clé en main ; pas d'observabilité ni d'interface
intégrées. \textbf{Manques :} RAGOps, traçabilité persistée, expérience
utilisateur complète.

\paragraph{OpenWebUI.} \textbf{Forces :} interface conversationnelle agréable,
auto-hébergeable. \textbf{Limites :} centrée sur le \emph{chat}, observabilité et
évaluation limitées. \textbf{Manques :} évaluation systématique de la qualité,
traçabilité fine, supervision de la santé du pipeline.

\paragraph{Haystack.} \textbf{Forces :} \emph{framework} RAG modulaire et
mature~\cite{haystack}. \textbf{Limites :} orienté développeur ; l'observabilité
et l'interface restent à construire. \textbf{Manques :} tour de contrôle
opérationnelle prête à l'emploi, visualisation pédagogique de l'exécution.

\paragraph{Plateformes de connaissance d'entreprise.} \textbf{Forces :}
intégration documentaire et gouvernance. \textbf{Limites :} recherche souvent
lexicale ; capacités génératives et de traçabilité du raisonnement limitées.
\textbf{Manques :} pipeline RAG agentique transparent et évalué.

\paragraph{Différenciation de \devpilot{}.} \devpilot{} se distingue en
\textbf{intégrant} l'application RAG \emph{et} son exploitation dans un système
unique, auto-hébergeable : base de connaissances privée par espace de travail,
réponses ancrées et citées, évaluation automatique, traçabilité complète
(RAGOps, explorateur et inspecteur de traces) et visualisation d'exécution
(AI Execution Studio). La transparence du pipeline de raisonnement est le c{\oe}ur
de la proposition de valeur.

\subsection{Étude comparative}
Le tableau~\ref{tab:comparaison} positionne \devpilot{} face aux solutions
analysées selon les capacités déterminantes pour une plateforme d'IA d'entreprise.

\begin{table}[H]
\centering\caption{Étude comparative des solutions}\label{tab:comparaison}
\scriptsize\renewcommand{\arraystretch}{1.25}
\begin{tabularx}{\textwidth}{@{}>{\raggedright\arraybackslash}p{3.1cm} *{6}{>{\centering\arraybackslash}X}@{}}
\toprule
\textbf{Capacité} & \textbf{ChatGPT} & \textbf{LangSmith} & \textbf{LlamaIndex} & \textbf{Haystack} & \textbf{OpenWebUI} & \textbf{\devpilot{}} \\
\midrule
Réponses ancrées sur vos documents & Limité & N/A & Oui & Oui & Partiel & \textbf{Oui} \\
Citations vérifiables & Non & N/A & Partiel & Partiel & Partiel & \textbf{Oui} \\
Traçabilité complète des requêtes & Non & Oui & Non & Non & Non & \textbf{Oui} \\
Isolation par espace de travail & Non & N/A & Non & Non & Partiel & \textbf{Oui} \\
Score d'hallucination & Non & Partiel & Non & Non & Non & \textbf{Oui} \\
Métriques de récupération & Non & Oui & Partiel & Partiel & Non & \textbf{Oui} \\
Chronologie d'exécution des agents & Non & Oui & Non & Non & Non & \textbf{Oui} \\
Auto-hébergement & Non & Partiel & Oui & Oui & Oui & \textbf{Oui} \\
Évaluation intégrée & Non & Partiel & Non & Partiel & Non & \textbf{Oui} \\
Observabilité opérationnelle (RAGOps) & Non & Partiel & Non & Non & Non & \textbf{Oui} \\
\bottomrule
\end{tabularx}
\end{table}

\subsection{Besoins fonctionnels}
Les exigences fonctionnelles (EF) sont synthétisées dans le
tableau~\ref{tab:ef}.

\begin{table}[H]
\centering\caption{Exigences fonctionnelles principales}\label{tab:ef}\small
\begin{tabularx}{\textwidth}{@{}l l X@{}}
\toprule
\textbf{Réf.} & \textbf{Domaine} & \textbf{Exigence} \\
\midrule
EF-1 & Authentification & Inscription, connexion et sessions sécurisées par jetons JWT \\
EF-2 & Autorisation & Contrôle d'accès par rôles (utilisateur, admin, super-admin) \\
EF-3 & Espaces & Création d'espaces isolés et gestion des membres \\
EF-4 & Documents & Téléversement, listing, suivi de statut et suppression \\
EF-5 & Ingestion & Traitement asynchrone : extraction, découpage, vectorisation, indexation \\
EF-6 & Récupération & Recherche sémantique et hybride sur la base vectorielle \\
EF-7 & Reclassement & Reclassement des candidats par pertinence (\emph{reranking}) \\
EF-8 & Génération & Production de réponses ancrées par un pipeline agentique \\
EF-9 & Citations & Génération de citations reliées aux passages utilisés \\
EF-10 & Conversation & Historique de conversations et de messages persistants \\
EF-11 & Streaming & Diffusion incrémentale des réponses (jeton par jeton) \\
EF-12 & Évaluation & Mesure de la fidélité, de la pertinence et du risque d'hallucination \\
EF-13 & Correction & Réparation automatique des réponses non fidèles \\
EF-14 & Administration & Tableaux de bord d'administration et journal d'audit \\
EF-15 & RAGOps & Supervision de la santé des pipelines et de la file de traitement \\
EF-16 & Traçabilité & Explorateur et inspecteur de traces de requêtes RAG \\
EF-17 & Visualisation & AI Execution Studio : visualisation animée de l'exécution \\
\bottomrule
\end{tabularx}
\end{table}

\subsection{Besoins non fonctionnels}
\begin{description}[leftmargin=2.4cm,style=nextline]
\item[Sécurité] Authentification JWT, autorisation par rôles, isolation stricte
des données par espace, journalisation d'audit et masquage systématique des
secrets dans les vues de débogage.
\item[Performance] API asynchrone non bloquante, déport du traitement lourd via
Celery, recherche par plus proches voisins approchés (HNSW) et \emph{streaming}
réduisant le \emph{time-to-first-token}.
\item[Scalabilité] \textit{Backend} sans état horizontalement extensible,
\emph{workers} Celery ajoutables, base vectorielle Qdrant partitionnable.
\item[Disponibilité] Conteneurisation, dépendances de santé explicites et
politiques de relance des tâches asynchrones.
\item[Maintenabilité] Architecture en couches, typage statique de bout en bout,
migrations de schéma versionnées (Alembic).
\item[Observabilité] Traçabilité complète, métriques Prometheus, tableaux de bord
Grafana, RAGOps et inspecteur de traces.
\item[Fiabilité] Repli de modèles de génération, agent correcteur et relance
automatique des ingestions échouées.
\item[Utilisabilité] Interface conversationnelle réactive, rendu Markdown, cartes
de citation interactives, thème clair/sombre.
\end{description}

\subsection{Contraintes}
\begin{description}[leftmargin=2.4cm,style=nextline]
\item[Techniques] Reproductibilité par conteneurs ; possibilité d'exécution
hors-ligne via des modèles locaux (Ollama) ; budget mémoire et calcul maîtrisé.
\item[Académiques] Démonstration d'une démarche d'ingénierie rigoureuse et d'un
raisonnement architectural.
\item[Temps] Réalisation dans le cadre temporel d'un projet de fin d'études.
\item[Ressources] Environnement d'exécution mono-machine, sans infrastructure
\emph{cloud} dédiée.
\end{description}

\subsection{Cahier des charges}
Le tableau~\ref{tab:cdc} relie les objectifs aux livrables vérifiables.

\begin{table}[H]
\centering\caption{Synthèse du cahier des charges}\label{tab:cdc}\small
\begin{tabularx}{\textwidth}{@{}>{\bfseries}l X@{}}
\toprule
Objectif & Livrable \\
\midrule
Ingestion & Pipeline asynchrone Celery (extraction, découpage, vectorisation, indexation Qdrant) \\
Récupération & Service de recherche hybride avec reclassement \\
Génération & Flux agentique de sept agents avec citations \\
Évaluation & Module d'évaluation (fidélité, pertinence, hallucination) \\
Exploitation & Tour de contrôle RAGOps \\
Traçabilité & Explorateur et inspecteur de traces \\
Visualisation & AI Execution Studio \\
Déploiement & Orchestration Docker Compose (8 services) \\
\bottomrule
\end{tabularx}
\end{table}

\subsection{Gestion des risques}
Le tableau~\ref{tab:risques} recense les principaux risques et leurs mesures de
mitigation.

\begin{table}[H]
\centering\caption{Gestion des risques}\label{tab:risques}\small
\begin{tabularx}{\textwidth}{@{}X c c X@{}}
\toprule
\textbf{Risque} & \textbf{Prob.} & \textbf{Impact} & \textbf{Mitigation} \\
\midrule
Hallucination du LLM & Moyenne & Élevé & Ancrage strict, évaluation de fidélité, agent correcteur \\
Indisponibilité du fournisseur LLM & Faible & Élevé & Modèle de repli (\emph{fallback}) ; embeddings locaux \\
Couverture d'\emph{embedding} incomplète & Moyenne & Moyen & Suivi RAGOps de la couverture, relance d'ingestion \\
Échec d'ingestion documentaire & Moyenne & Moyen & Relance Celery (3 tentatives), centre des échecs \\
Volume de \emph{node\_modules} conteneur & Faible & Faible & Procédure de réinstallation documentée \\
Dégradation des performances & Faible & Moyen & Recherche HNSW, \emph{streaming}, chargement différé \\
\bottomrule
\end{tabularx}
\end{table}

\subsection{Conclusion}
Ce chapitre a établi que le besoin n'est pas un \emph{chatbot} supplémentaire,
mais une plateforme d'IA \emph{fondée}, \emph{mesurable} et \emph{observable}.
L'étude comparative montre que \devpilot{} occupe une position distinctive en
unifiant moteur RAG, exploitation et traçabilité. Les exigences non fonctionnelles
— sécurité, scalabilité, observabilité — sont aussi structurantes que les
exigences fonctionnelles. Le chapitre suivant traduit ces besoins en modèles de
conception.

\chapter{Analyse et conception}
\label{chap:conception}

Ce chapitre formalise la conception de \devpilot{} : identification des acteurs,
cas d'utilisation, modèle du domaine, diagrammes de classes, de séquence, modèle
de données, et vues d'architecture (logique, physique, déploiement, composants,
paquets, communication). Une attention particulière est portée à la modélisation
des quatre modules distinctifs : le moteur RAG, RAGOps, le système de traçabilité
et l'AI Execution Studio. Chaque diagramme est introduit par son explication, ses
hypothèses de modélisation, son code PlantUML et son interprétation.

\section{Identification des acteurs}
Trois acteurs humains interagissent avec le système, selon une hiérarchie de
droits : l'\textbf{Utilisateur} (membre d'une équipe d'ingénierie), l'\textbf{Administrateur}
(gestion de la plateforme) et le \textbf{Super-administrateur} (observabilité
approfondie). Trois acteurs systèmes externes participent au pipeline :
\textbf{Qdrant} (base vectorielle), \textbf{Ollama} (\emph{embeddings}) et
\textbf{Gemini} (génération).

\section{Cas d'utilisation}
Les cas d'utilisation clés sont : gérer un espace de travail, téléverser et suivre
un document, poser une question (RAG) et recevoir une réponse citée en
\emph{streaming}, consulter l'historique, superviser les pipelines (RAGOps) et
inspecter une trace. Le tableau~\ref{tab:uc} détaille le cas pivot.

\begin{table}[H]
\centering\caption{Cas d'utilisation détaillé : « Poser une question (RAG) »}
\label{tab:uc}\small
\begin{tabularx}{\textwidth}{@{}>{\bfseries}l X@{}}
\toprule
Acteur & Utilisateur authentifié, membre d'un espace de travail \\
Pré-condition & L'espace contient au moins un document indexé \\
Déclencheur & Soumission d'une question dans l'interface de chat \\
Scénario nominal & Router, réécriture, récupération, reclassement, génération, évaluation, correction, diffusion en \emph{streaming} \\
Post-condition & Message, agents, passages et évaluation persistés \\
Exceptions & Aucun passage pertinent $\rightarrow$ repli explicite ; échec LLM $\rightarrow$ modèle de \emph{fallback} \\
\bottomrule
\end{tabularx}
\end{table}

\section{Diagramme de cas d'utilisation}
\paragraph{Explication.} Ce diagramme présente les acteurs et leurs interactions
avec les fonctionnalités principales, ainsi que la relation d'héritage entre rôles.
\paragraph{Hypothèses.} Les acteurs systèmes (Qdrant, Ollama, Gemini) sont
représentés comme acteurs secondaires sollicités par le cas « Poser une question ».
Le super-administrateur hérite de l'administrateur, qui hérite de l'utilisateur.

\begin{lstlisting}[style=plantuml,caption={Diagramme de cas d'utilisation},label={lst:uc}]
@startuml
left to right direction
actor Utilisateur as U
actor Administrateur as A
actor "Super-administrateur" as S
actor Qdrant as Q
actor Gemini as G

A --|> U
S --|> A

rectangle "DevPilot AI" {
  usecase "Gérer un espace" as UC1
  usecase "Téléverser un document" as UC2
  usecase "Poser une question (RAG)" as UC3
  usecase "Consulter l'historique" as UC4
  usecase "Administrer les utilisateurs" as UC5
  usecase "Superviser (RAGOps)" as UC6
  usecase "Inspecter une trace" as UC7
}

U --> UC1
U --> UC2
U --> UC3
U --> UC4
A --> UC5
A --> UC6
S --> UC7
UC3 ..> Q : <<include>>
UC3 ..> G : <<include>>
@enduml
\end{lstlisting}

\paragraph{Interprétation.} L'utilisateur dispose des cas métier (espaces,
documents, RAG, historique). L'administrateur ajoute la supervision, le
super-administrateur l'inspection fine. Le cas RAG inclut systématiquement la
recherche vectorielle et la génération.
\plantumlfig{Diagramme de cas d'utilisation global}{fig:uc}

\section{Modèle du domaine}
\paragraph{Explication.} Le modèle du domaine articule les concepts métier
centraux et leurs relations, indépendamment de la persistance.
\paragraph{Hypothèses.} La \emph{Trace} est élevée au rang de concept du domaine
(agrégat des exécutions d'agents, passages récupérés et évaluation), choix
déterminant pour l'observabilité.

\begin{lstlisting}[style=plantuml,caption={Modèle du domaine},label={lst:domaine}]
@startuml
class Utilisateur
class EspaceTravail
class Adhesion
class Document
class Chunk
class Conversation
class Message
class Trace

Utilisateur "1" -- "0..*" Adhesion
EspaceTravail "1" -- "0..*" Adhesion
EspaceTravail "1" -- "0..*" Document
Document "1" *-- "0..*" Chunk
Utilisateur "1" -- "0..*" Conversation
Conversation "1" *-- "1..*" Message
Message "1" -- "0..1" Trace
@enduml
\end{lstlisting}

\paragraph{Interprétation.} Un utilisateur accède aux espaces via des adhésions ;
un espace contient des documents découpés en \emph{chunks} ; les conversations
agrègent des messages, et chaque message assistant porte une trace.
\plantumlfig{Modèle du domaine}{fig:domaine}

\section{Diagrammes de classes}
\subsection{Entités persistantes}
\paragraph{Explication.} Ce diagramme présente les onze entités persistantes avec
leurs attributs clés et leurs cardinalités.
\paragraph{Hypothèses.} Les identifiants sont des UUID ; les attributs techniques
communs (horodatages) sont omis pour la lisibilité.

\begin{lstlisting}[style=plantuml,caption={Diagramme de classes — entités persistantes},label={lst:classes}]
@startuml
class users { +id : UUID\n+email\n+role\n+is_active }
class workspaces { +id : UUID\n+name\n+owner_id }
class workspace_members { +id : UUID\n+role }
class documents { +id : UUID\n+filename\n+status }
class chunks { +id : UUID\n+chunk_index\n+vector_id }
class conversations { +id : UUID\n+title }
class messages { +id : UUID\n+role\n+content }
class agent_runs { +id : UUID\n+agent_type\n+latency_ms }
class retrieved_chunks { +id : UUID\n+rank\n+score }
class evaluations { +id : UUID\n+faithfulness\n+hallucination_score }
class llm_usage { +id : UUID\n+total_tokens\n+cost_usd }

users "1" -- "0..*" workspaces
users "1" -- "0..*" workspace_members
workspaces "1" -- "0..*" workspace_members
workspaces "1" *-- "0..*" documents
documents "1" *-- "0..*" chunks
users "1" -- "0..*" conversations
conversations "1" *-- "1..*" messages
messages "1" -- "0..*" agent_runs
messages "1" -- "0..*" retrieved_chunks
messages "1" -- "1" evaluations
messages "1" -- "0..*" llm_usage
@enduml
\end{lstlisting}

\paragraph{Interprétation.} Les compositions (\emph{chunks} dans \emph{documents},
\emph{messages} dans \emph{conversations}) traduisent un cycle de vie lié ; la
relation 1--1 entre \texttt{messages} et \texttt{evaluations} formalise l'unicité
de l'évaluation par réponse.
\plantumlfig{Diagramme de classes des entités persistantes}{fig:classes}

\subsection{Moteur agentique}
\paragraph{Explication.} Ce diagramme modélise le flux agentique : un orchestrateur
coordonne sept agents, appuyés par des services applicatifs.
\paragraph{Hypothèses.} Les agents partagent une interface commune d'exécution ;
seuls les services principaux sont représentés.

\begin{lstlisting}[style=plantuml,caption={Diagramme de classes — moteur agentique},label={lst:agents}]
@startuml
class AgenticRAGWorkflow {
  +run(question) : Result
}
interface Agent {
  +execute(ctx)
}
class RouterAgent
class QueryRewriterAgent
class RetrievalAgent
class RerankerAgent
class GeneratorAgent
class EvaluatorAgent
class CorrectorAgent
class RetrievalService
class RerankingService
class EvaluationService
class RagTracesService

RouterAgent ..|> Agent
QueryRewriterAgent ..|> Agent
RetrievalAgent ..|> Agent
RerankerAgent ..|> Agent
GeneratorAgent ..|> Agent
EvaluatorAgent ..|> Agent
CorrectorAgent ..|> Agent

AgenticRAGWorkflow o-- RouterAgent
AgenticRAGWorkflow o-- QueryRewriterAgent
AgenticRAGWorkflow o-- RetrievalAgent
AgenticRAGWorkflow o-- RerankerAgent
AgenticRAGWorkflow o-- GeneratorAgent
AgenticRAGWorkflow o-- EvaluatorAgent
AgenticRAGWorkflow o-- CorrectorAgent
RetrievalAgent --> RetrievalService
RerankerAgent --> RerankingService
EvaluatorAgent --> EvaluationService
AgenticRAGWorkflow --> RagTracesService
@enduml
\end{lstlisting}

\paragraph{Interprétation.} L'orchestrateur agrège les sept agents (agrégation
\texttt{o--}) qui réalisent une interface commune ; les agents délèguent aux
services métier la récupération, le reclassement et l'évaluation.
\plantumlfig{Diagramme de classes du moteur agentique}{fig:agents}

\section{Diagrammes de séquence}
\subsection{Ingestion documentaire asynchrone}
\paragraph{Explication.} Séquence du téléversement à l'indexation, exécutée hors
du chemin critique par Celery.
\paragraph{Hypothèses.} Le stockage de fichiers et la base relationnelle sont
distingués ; la réponse HTTP est immédiate (202).

\begin{lstlisting}[style=plantuml,caption={Séquence — ingestion documentaire},label={lst:seq-ingest}]
@startuml
actor Client
participant "API (FastAPI)" as API
participant "Stockage" as ST
participant "Celery" as CEL
participant "Qdrant" as QD

Client -> API : POST /documents (fichier)
activate API
API -> ST : persister fichier + métadonnées
API -> CEL : enqueue process_document_task
API --> Client : 202 Accepted (statut=pending)
deactivate API
activate CEL
CEL -> ST : lire + extraire le texte
CEL -> CEL : découper + vectoriser (768-d)
CEL -> QD : upsert vecteurs
CEL -> ST : statut = indexed
deactivate CEL
@enduml
\end{lstlisting}

\paragraph{Interprétation.} L'API rend la main immédiatement ; le \emph{worker}
réalise l'extraction, le découpage, la vectorisation et l'indexation, puis met à
jour le statut du document.
\plantumlfig{Diagramme de séquence — ingestion asynchrone}{fig:seq-ingest}

\subsection{Requête RAG en streaming}
\paragraph{Explication.} Séquence de la requête RAG diffusée en \emph{streaming}
SSE à travers le flux agentique.
\paragraph{Hypothèses.} Les sept agents sont synthétisés au niveau du
\texttt{Workflow} ; les événements SSE sont émis au fil de l'eau.

\begin{lstlisting}[style=plantuml,caption={Séquence — requête RAG en streaming},label={lst:seq-rag}]
@startuml
actor Client
participant "API (SSE)" as API
participant "Workflow" as WF
participant "Qdrant" as QD
participant "Gemini" as LLM

Client -> API : POST /chat/stream (question)
activate API
API -> WF : exécuter le pipeline
activate WF
WF -> WF : router + réécriture
WF -> QD : récupérer candidats (k=15)
QD --> WF : passages + scores
WF -> WF : reclasser (top 5)
API --> Client : event: trace / citations
WF -> LLM : générer (stream)
LLM --> WF : flux de jetons
API --> Client : event: token (xN)
WF -> WF : évaluer + corriger
API --> Client : event: evaluation / message / done
deactivate WF
deactivate API
@enduml
\end{lstlisting}

\paragraph{Interprétation.} La diffusion entrelace les étapes du pipeline et les
événements SSE : les citations précèdent les jetons, suivis de l'évaluation et de
l'événement de clôture.
\plantumlfig{Diagramme de séquence — requête RAG en streaming}{fig:seq-rag}

\subsection{Inspection d'une trace}
\paragraph{Explication.} Séquence de consultation d'une trace par
l'administrateur.
\paragraph{Hypothèses.} L'inspecteur agrège quatre appels parallèles (trace,
récupération, reclassement, évaluation).

\begin{lstlisting}[style=plantuml,caption={Séquence — inspection de trace},label={lst:seq-trace}]
@startuml
actor Admin
participant "Frontend" as FE
participant "API Admin" as API
database "PostgreSQL" as PG

Admin -> FE : ouvrir /admin/rag-traces/{id}
FE -> API : GET trace + retrieval + reranking + evaluation
activate API
API -> PG : lire la lignée du message
PG --> API : agent_runs, retrieved_chunks, evaluations
API --> FE : trace complète
deactivate API
FE --> Admin : reconstruction des 9 sections
@enduml
\end{lstlisting}

\paragraph{Interprétation.} La trace est reconstruite intégralement à partir des
données persistées, sans dépendance à des journaux externes.
\plantumlfig{Diagramme de séquence — inspection de trace}{fig:seq-trace}

\section{Modèle de données}
\paragraph{Explication.} Le modèle relationnel comprend onze tables ; la
figure~\ref{fig:erd} en présente la vue entité-association.
\paragraph{Hypothèses.} Les clés étrangères sont implicites dans les associations ;
les colonnes techniques sont omises.

\begin{lstlisting}[style=plantuml,caption={Modèle de données (ERD)},label={lst:erd}]
@startuml
entity users
entity workspaces
entity workspace_members
entity documents
entity chunks
entity conversations
entity messages
entity agent_runs
entity retrieved_chunks
entity evaluations
entity llm_usage

users ||--o{ workspaces
users ||--o{ workspace_members
workspaces ||--o{ workspace_members
workspaces ||--o{ documents
documents ||--o{ chunks
users ||--o{ conversations
conversations ||--o{ messages
messages ||--o{ agent_runs
messages ||--o{ retrieved_chunks
messages ||--|| evaluations
messages ||--o{ llm_usage
@enduml
\end{lstlisting}

\paragraph{Interprétation.} Le schéma matérialise la chaîne de traçabilité
\texttt{message} $\rightarrow$ \texttt{agent\_runs} $\rightarrow$
\texttt{retrieved\_chunks} $\rightarrow$ \texttt{evaluations}, complétée par
\texttt{llm\_usage}.
\plantumlfig{Modèle de données — vue entité-association}{fig:erd}

\section{Architecture logique}
L'application s'organise en couches aux responsabilités strictes : présentation
(\emph{frontend}), application (API et services), domaine IA (pipeline agentique),
traitement asynchrone (Celery) et données (PostgreSQL, Qdrant, stockage). Les
dépendances ne circulent que du haut vers le bas, condition de la maintenabilité.

\section{Architecture physique}
Les composants sont déployés sous forme de conteneurs indépendants communiquant
par le réseau : services applicatifs sans état, magasins de données persistants et
services d'IA (Ollama local, Gemini distant). Cette séparation autorise une montée
en charge ciblée de chaque composant.

\section{Architecture de déploiement}
\paragraph{Explication.} Diagramme de déploiement des huit services conteneurisés,
de l'hôte Ollama et de l'API Gemini externe.
\paragraph{Hypothèses.} Un n{\oe}ud unique héberge les conteneurs ; Ollama
s'exécute sur l'hôte, Gemini est un service \emph{cloud}.

\begin{lstlisting}[style=eraser,caption={Diagramme de déploiement (Eraser)},label={lst:deploy}]
colorMode bold
styleMode shadow
direction down

Docker Host [icon: docker, color: blue] {
  frontend [icon: nextjs, label: "frontend (3001)"]
  backend [icon: python, label: "backend uvicorn (8000)"]
  celery_worker [icon: python]
  postgres [icon: postgresql, label: "postgres:16"]
  redis [icon: redis, label: "redis:7"]
  qdrant [icon: database, label: "qdrant 1.12"]
  prometheus [icon: prometheus]
  grafana [icon: grafana]
}
Ollama Host [icon: server, color: teal] { ollama [icon: cpu, label: "Ollama :11434"] }
Gemini Cloud [icon: google, color: orange] { gemini [icon: cloud, label: "Gemini API"] }

frontend > backend: HTTPS / SSE
backend > postgres
backend > redis
backend > qdrant
backend > ollama
backend > gemini
celery_worker > redis
celery_worker > postgres
celery_worker > qdrant
backend > prometheus
prometheus > grafana
\end{lstlisting}

\paragraph{Interprétation.} Le \emph{backend} et le \emph{worker} partagent les
magasins de données ; Redis relie l'API au \emph{worker} ; l'observabilité repose
sur Prometheus et Grafana.
\eraserfig{Diagramme de déploiement}{fig:deploy}

\section{Diagramme de composants}
\paragraph{Explication.} Vue des composants applicatifs et de leurs dépendances.
\paragraph{Hypothèses.} Les interfaces réseau (REST, SSE) sont représentées par
les flèches de dépendance.

\begin{lstlisting}[style=eraser,caption={Diagramme de composants (Eraser)},label={lst:comp}]
colorMode bold
styleMode shadow
direction down

Frontend [icon: nextjs, color: blue]
API [icon: python, color: blue, label: "API FastAPI"]
Workflow [icon: cpu, color: purple, label: "Pipeline agentique"]
Celery [icon: refresh-cw, color: orange]
Redis [icon: redis, color: orange]
Postgres [icon: postgresql, color: green]
Qdrant [icon: database, color: green]
Ollama [icon: server, color: teal]
Gemini [icon: google, color: teal]
Observability [icon: activity, color: red, label: "Prometheus / Grafana"]

Frontend > API
API > Workflow
API > Postgres
API > Redis
Redis > Celery
Workflow > Qdrant
Workflow > Ollama
Workflow > Gemini
Celery > Postgres
Celery > Qdrant
API > Observability
\end{lstlisting}

\paragraph{Interprétation.} Le \emph{frontend} ne dépend que de l'API ; le
pipeline agentique concentre les dépendances vers Qdrant et les fournisseurs d'IA.
\eraserfig{Diagramme de composants}{fig:comp}

\section{Diagramme de paquets}
\paragraph{Explication.} Organisation modulaire du code \emph{backend} et
\emph{frontend}.
\paragraph{Hypothèses.} Seuls les paquets principaux sont représentés.

\begin{lstlisting}[style=plantuml,caption={Diagramme de paquets},label={lst:pkg}]
@startuml
package "backend.app" {
  package api
  package core
  package db
  package models
  package schemas
  package services
  package agents
  package workers
  package utils
}
package "frontend" {
  package app
  package components
  package features
  package lib
}
api --> services
services --> models
services --> agents
agents --> utils
workers --> services
app --> features
features --> components
features --> lib
@enduml
\end{lstlisting}

\paragraph{Interprétation.} Le découpage en paquets reflète la séparation des
responsabilités : routes, logique métier, agents et accès aux données côté
serveur ; pages, modules fonctionnels, composants et bibliothèques côté client.
\plantumlfig{Diagramme de paquets}{fig:pkg}

\section{Diagramme de communication}
\paragraph{Explication.} Vue de collaboration numérotée de la requête RAG.
\paragraph{Hypothèses.} Les messages sont ordonnés selon la séquence
d'exécution.

\begin{lstlisting}[style=plantuml,caption={Diagramme de communication — requête RAG},label={lst:comm}]
@startuml
object Client
object API
object Workflow
object Qdrant
object Gemini

Client -> API : 1: question
API -> Workflow : 2: exécuter
Workflow -> Qdrant : 3: récupérer
Qdrant -> Workflow : 4: passages
Workflow -> Gemini : 5: générer
Gemini -> Workflow : 6: jetons
Workflow -> API : 7: trace + réponse
API -> Client : 8: stream + citations
@enduml
\end{lstlisting}

\paragraph{Interprétation.} La numérotation explicite l'ordre des échanges et met
en évidence le rôle central de l'orchestrateur dans la coordination.
\plantumlfig{Diagramme de communication}{fig:comm}

\section{Conclusion}
La conception traduit les besoins en modèles exploitables : cas d'utilisation,
modèle du domaine, diagrammes de classes (entités et moteur agentique), scénarios
d'interaction, modèle de données à onze tables et vues d'architecture. La
modélisation explicite de la trace et du flux agentique prépare directement
l'architecture détaillée au chapitre~\ref{chap:archi}.

\chapter{Architecture et réalisation}
\label{chap:archi}

Ce chapitre constitue le c{\oe}ur du travail d'ingénierie. Il détaille
l'architecture globale et par couche, les pipelines RAG, RAGOps, de traçabilité
et de l'AI Execution Studio, les flux d'exécution, ainsi que les aspects
transverses de sécurité, de performance, de scalabilité et d'observabilité.

\section{Architecture globale}
\devpilot{} est un système orienté services déployé par conteneurs. La vision
architecturale repose sur quatre principes : séparation des responsabilités en
couches, asynchronisme pour découpler le chemin critique du traitement lourd,
traçabilité native (chaque étape produit une trace persistée) et neutralité
vis-à-vis des fournisseurs d'IA.

\paragraph{Explication.} Le diagramme suivant présente le flux de bout en bout à
travers les composants principaux.
\paragraph{Hypothèses.} Les magasins de données sont distingués ; le pipeline
agentique est synthétisé en un composant.

\begin{lstlisting}[style=eraser,caption={Architecture globale et flux de données (Eraser)},label={lst:global}]
colorMode bold
styleMode shadow
direction down

User [icon: user]
Client [icon: monitor, color: blue] { Frontend [icon: nextjs] }
API Gateway [icon: server, color: blue] { FastAPI [icon: python]; Auth [icon: lock] }
Orchestration [icon: workflow, color: purple] { Workflow [icon: cpu, label: "7 agents"] }
Async [icon: refresh-cw, color: orange] { Redis [icon: redis]; Celery [icon: python] }
Data [icon: database, color: green] { Postgres [icon: postgresql]; Qdrant [icon: database] }
AI [icon: cpu, color: teal] { Ollama [icon: server]; Gemini [icon: google] }
Obs [icon: activity, color: red] { Prometheus [icon: prometheus]; Grafana [icon: grafana] }

User > Frontend: question
Frontend > FastAPI: HTTPS / SSE
FastAPI > Workflow
FastAPI > Redis: enqueue
Redis > Celery
Workflow > Qdrant: recherche
Workflow > Ollama: embeddings
Workflow > Gemini: génération
Celery > Qdrant: upsert
FastAPI > Postgres
Celery > Postgres
FastAPI > Prometheus
Prometheus > Grafana
\end{lstlisting}

\paragraph{Interprétation.} La requête synchrone traverse l'API et le pipeline ;
l'ingestion documentaire est déléguée à Celery via Redis. PostgreSQL conserve
l'état et les traces, Qdrant les vecteurs.
\eraserfig{Architecture système globale}{fig:global}

\section{Architecture Backend}
Le \emph{backend} FastAPI s'organise en couches aux frontières nettes :
\texttt{api} (routes et intégration HTTP), \texttt{core} (configuration, sécurité,
exceptions), \texttt{db} (session et base ORM), \texttt{models} (modèles
SQLAlchemy), \texttt{schemas} (schémas Pydantic), \texttt{services} (logique
applicative), \texttt{agents} (orchestration agentique), \texttt{workers} (tâches
Celery) et \texttt{utils}. Les routes restent minces et délèguent aux services, ce
qui isole la logique métier et facilite les tests.

\section{Architecture Frontend}
Le \emph{frontend} Next.js (\emph{App Router}, React~19) adopte une organisation
par modules fonctionnels (\texttt{features/}), un système de design fondé sur des
jetons CSS (« Horizon », thème clair/sombre) et la gestion de l'état serveur par
React Query. Les visualisations lourdes (AI Execution Studio, Architecture
Explorer) sont chargées de manière différée (\emph{lazy loading}) afin de
préserver le poids initial des routes.

\section{Architecture Base de données}
Le modèle relationnel comprend onze tables gérées par SQLAlchemy et versionnées
par Alembic. \devpilot{} repose sur un \textbf{stockage dual} : PostgreSQL
conserve le texte canonique des fragments et une référence \texttt{vector\_id}
vers le point Qdrant correspondant, tandis que Qdrant stocke les vecteurs
indexés en HNSW (distance cosinus). Cette séparation préserve l'intégrité
référentielle tout en confiant la recherche vectorielle à un moteur spécialisé.

\section{Architecture API}
L'API est versionnée sous \texttt{/api/v1} et organisée en routeurs
(authentification, espaces, documents, chat, système, administration, RAGOps,
traces). Elle expose une interface REST pour les opérations CRUD et une interface
\emph{streaming} (SSE) pour la génération. L'autorisation repose sur des
dépendances FastAPI appliquant le contrôle d'accès par rôles.

\section{Architecture RAG}
La génération suit un pipeline en plusieurs étapes : \textbf{ingestion} (extraction
du texte), \textbf{découpage} en fragments à fenêtres chevauchantes,
\textbf{vectorisation} (\emph{embeddings} \texttt{nomic-embed-text}, 768
dimensions), \textbf{indexation} dans Qdrant, \textbf{récupération} hybride
(sémantique et lexicale) d'un ensemble large de candidats, \textbf{reclassement}
fin, \textbf{construction du prompt} (système + contexte + question) et
\textbf{génération} ancrée et citée. La récupération suit une stratégie « rappel
large puis précision fine » : quinze candidats sont remontés puis réduits aux cinq
meilleurs.

\paragraph{Explication.} Le diagramme suivant détaille le flux agentique de sept
agents coopérants.
\paragraph{Hypothèses.} Chaque agent est isolé et persisté ; le correcteur
n'intervient qu'en cas de réponse non fidèle.

\begin{lstlisting}[style=eraser,caption={Pipeline agentique RAG (Eraser)},label={lst:pipe}]
colorMode bold
styleMode shadow
direction right

Question [icon: message-circle, color: blue]
Pipeline [icon: workflow, color: purple] {
  Router [icon: git-branch]; QueryRewriter [icon: edit]; Retrieval [icon: search]
  Reranker [icon: bar-chart]; Generator [icon: sparkles]; Evaluator [icon: shield]; Corrector [icon: tool]
}
Qdrant [icon: database, color: green, label: "k=15"]
Gemini [icon: google, color: teal]
Answer [icon: check-circle, color: green, label: "ancrée + citée"]

Question > Router > QueryRewriter > Retrieval
Retrieval > Qdrant: 15 candidats
Retrieval > Reranker > Generator
Generator > Gemini: stream
Generator > Evaluator
Evaluator > Corrector: si non fidèle
Evaluator > Answer: si fidèle
Corrector > Answer
\end{lstlisting}

\paragraph{Interprétation.} Le flux est linéaire jusqu'à l'évaluation ; le
correcteur constitue une boucle de réparation conditionnelle garantissant la
fidélité de la réponse diffusée.
\eraserfig{Pipeline de génération agentique}{fig:pipe}

\section{Architecture RAGOps}
RAGOps constitue la couche d'\emph{exploitation} du pipeline RAG — le pont entre
architecture et observabilité. Là où un tableau de bord classique rapporte des
métriques métier, RAGOps rapporte la santé et la qualité du système d'IA lui-même.
Il agrège plusieurs dimensions :

\begin{description}[leftmargin=2.6cm,style=nextline]
\item[Health Monitoring] Santé de la collection Qdrant et état des services
(\texttt{/admin/ragops/qdrant/health}, \texttt{ops-health}, \texttt{vector-health}).
\item[Retrieval Analytics] Stratégie de récupération, nombre de candidats et
qualité du contexte récupéré.
\item[Evaluation Intelligence] Fidélité moyenne, comptage des réponses à risque
d'hallucination et des réponses corrigées.
\item[Embedding Analytics] Couverture d'\emph{embedding} par espace
(\texttt{chunks\_with\_vector\_id} / \texttt{total\_chunks}).
\item[Workspace Health] Score de santé RAG par espace (0--100) combinant ratio
d'indexation, couverture, fidélité et fiabilité.
\item[Agent Monitoring] Latence moyenne par agent et comptabilité des jetons et
coûts.
\item[Failure Center] Documents en échec et relance de l'ingestion
(\texttt{POST /admin/ragops/documents/\{id\}/retry}).
\item[Observability] Métriques Prometheus et tableaux de bord Grafana.
\end{description}

RAGOps est essentiel dans les systèmes d'IA d'entreprise car la qualité d'un
pipeline RAG dérive silencieusement (documents non indexés, couverture partielle,
hausse des hallucinations) : sans observabilité spécifique à l'IA, ces régressions
restent invisibles jusqu'à l'incident.

\section{Architecture Traceability}
La traçabilité est modélisée comme une chaîne dirigée centrée sur le message
assistant.

\paragraph{Explication.} Le diagramme présente la lignée d'exécution persistée.
\paragraph{Hypothèses.} Chaque message assistant agrège ses exécutions d'agents,
ses passages, son évaluation et sa comptabilité d'usage.

\begin{lstlisting}[style=eraser,caption={Chaîne de traçabilité (Eraser)},label={lst:trace}]
colorMode bold
styleMode shadow
direction right

Message [icon: message-square, color: blue]
Lineage [icon: layers, color: purple] {
  AgentRuns [icon: cpu, label: "agent_runs"]
  RetrievedChunks [icon: file-text, label: "retrieved_chunks"]
  Evaluations [icon: shield, label: "evaluations"]
  LLMUsage [icon: dollar-sign, label: "llm_usage"]
}
Storage [icon: database, color: green, label: "PostgreSQL"]
Inspector [icon: search, color: teal, label: "Trace Inspector"]

Message > AgentRuns: 1..*
Message > RetrievedChunks: 1..*
Message > Evaluations: 1..1
Message > LLMUsage: 0..*
AgentRuns > Storage
Evaluations > Storage
Storage > Inspector: reconstruction
\end{lstlisting}

\paragraph{Interprétation.} Toute réponse est \emph{reconstructible a posteriori}
à partir de ces données, sans dépendance à des journaux externes. Cela fonde
quatre propriétés : \textbf{auditabilité} (justifier une réponse), \textbf{reproductibilité}
(rejouer et comparer), \textbf{explicabilité} (exposer le raisonnement) et
\textbf{débogage/gouvernance} (diagnostiquer les régressions et tracer les actions
d'administration via le journal d'audit).
\eraserfig{Chaîne de traçabilité}{fig:trace}

\section{Architecture AI Execution Studio}
L'AI Execution Studio transforme le pipeline RAG en une expérience animée en temps
réel : au lieu d'afficher des métriques, il donne à voir le raisonnement du système.
Il est piloté par les \emph{mêmes API réelles} que le chat (flux SSE + endpoints de
trace). Il visualise séquentiellement les dix étapes : question, \emph{embedding},
recherche vectorielle, récupération des fragments, reclassement, construction du
contexte, construction du prompt, génération (\emph{streaming} de jetons),
évaluation (jauges animées) et persistance de la trace. Sa valeur est double :
\textbf{pédagogique} (comprendre le RAG en une trentaine de secondes) et
\textbf{d'observabilité} (visualiser où se dépensent la latence et la qualité).

\section{Flux d'ingestion}
L'ingestion est entièrement asynchrone (cf. figure~\ref{fig:seq-ingest}) : l'API
persiste le document avec le statut \texttt{pending} et délègue à Celery la tâche
\texttt{process\_document\_task}, relancée jusqu'à trois fois en cas d'échec, avant
passage au statut \texttt{failed}.

\section{Flux de recherche}
La recherche vectorise la requête réécrite, interroge Qdrant pour un ensemble de
candidats, fusionne signaux sémantiques et lexicaux (récupération hybride), puis
reclasse pour ne conserver que les passages les plus pertinents dans la limite du
budget de contexte (6000 caractères).

\section{Flux de génération}
Le générateur construit le prompt final (système + contexte + question) et produit
une réponse ancrée. Les fragments retenus sont reliés à des marqueurs de citation
\texttt{[n]} ; un mécanisme garantit la présence d'au moins une citation lorsque
du contexte est disponible.

\section{Flux de streaming SSE}
La diffusion repose sur \emph{Server-Sent Events} : le \emph{backend} renvoie une
\texttt{StreamingResponse} (\texttt{text/event-stream}, en-tête
\texttt{X-Accel-Buffering: no}) émettant une séquence d'événements typés
(\texttt{start}, \texttt{trace}, \texttt{token}, \texttt{citations},
\texttt{evaluation}, \texttt{correction}, \texttt{message}, \texttt{done},
\texttt{error}). Le \emph{streaming} repose sur l'API \texttt{streamGenerateContent}
de Gemini ; côté client, le flux est consommé via \texttt{fetch} et un
\texttt{ReadableStream}. Le listing~\ref{lst:sse} illustre le générateur côté
serveur.

\begin{lstlisting}[style=code,language=Python,caption={Génération du flux SSE (extrait)},label={lst:sse}]
async def event_stream() -> AsyncIterator[str]:
    async for event, data in chat_service.stream_chat_events(...):
        yield sse_event(event, data)

return StreamingResponse(
    event_stream(),
    media_type="text/event-stream",
    headers={"X-Accel-Buffering": "no"},
)
\end{lstlisting}

\section{Sécurité}
La sécurité est multicouche : authentification JWT, autorisation par rôles,
isolation des données par espace de travail (toute requête est cantonnée aux
espaces dont l'utilisateur est membre), journalisation des actions
d'administration et masquage systématique des secrets (clés contenant
\emph{password}, \emph{secret}, \emph{token}, \emph{api\_key}) dans les vues de
débogage et de traçabilité.

\section{Authentification JWT}
L'authentification repose sur des jetons JWT signés, émis à la connexion et
transmis dans l'en-tête \texttt{Authorization: Bearer}. Les dépendances FastAPI
résolvent l'utilisateur courant et vérifient son activité avant tout accès aux
ressources protégées.

\section{Gestion des permissions}
Trois rôles structurent l'accès : \texttt{user}, \texttt{admin} et
\texttt{super\_admin}. Le contrôle s'applique à deux niveaux : par route (les
\emph{endpoints} d'administration, RAGOps et traces exigent un rôle privilégié) et
par espace de travail (l'appartenance conditionne l'accès aux documents et
conversations).

\section{Performance}
Les leviers de performance sont l'asynchronisme bout-en-bout (E/S non bloquantes,
déport du calcul lourd vers Celery), la recherche HNSW sous-linéaire, le
\emph{streaming} (réduction du \emph{time-to-first-token}), ainsi que le
chargement différé et le découpage des \emph{bundles} côté \emph{frontend}
(routes de 110 à 175~kB environ).

\section{Scalabilité}
Le \emph{backend} sans état est horizontalement extensible derrière un répartiteur
de charge ; les \emph{workers} Celery se multiplient pour absorber les pics
d'ingestion ; Qdrant supporte le partitionnement et PostgreSQL la réplication en
lecture.

\section{Observabilité}
L'observabilité opère à trois niveaux complémentaires : métriques système
(Prometheus, tableaux Grafana), exploitation applicative (RAGOps) et traçabilité
des requêtes (explorateur et inspecteur de traces), auxquels s'ajoute la
visualisation pédagogique de l'AI Execution Studio.

\section{Conclusion}
L'architecture de \devpilot{} traduit ses principes directeurs en un système
concret, en couches, asynchrone, traçable et neutre vis-à-vis des fournisseurs
d'IA. Les pipelines documentaire, de récupération, de génération agentique et
d'évaluation, instrumentés par RAGOps et la traçabilité, forment un tout cohérent.
Le chapitre suivant justifie les choix technologiques et détaille l'implémentation.

\chapter{Technologies et implémentation}
\label{chap:techno}

Ce chapitre justifie les choix technologiques selon une grille uniforme
(pourquoi, comment, bénéfices, limites, alternatives) puis détaille les points
saillants de l'implémentation.

\section{Présentation de la pile technologique}
Le tableau~\ref{tab:stack} synthétise la pile par niveau.

\begin{table}[H]
\centering\caption{Pile technologique de DevPilot AI}\label{tab:stack}\small
\begin{tabularx}{\textwidth}{@{}>{\bfseries}l l X@{}}
\toprule
Niveau & Technologie & Rôle \\
\midrule
Backend & FastAPI, Pydantic & API REST et \emph{streaming}, validation \\
 & SQLAlchemy, Alembic & ORM asynchrone et migrations \\
Base de données & PostgreSQL 16 & État relationnel et traces \\
Messagerie & Redis 7 & Courtier Celery et \emph{cache} \\
Asynchrone & Celery & Traitement documentaire \\
Vectoriel & Qdrant 1.12 & Recherche HNSW (cosinus) \\
Modèles IA & Ollama, Gemini & \emph{Embeddings} locaux, génération \\
Frontend & Next.js 15, React 19, TypeScript & Interface, chat, tableaux de bord \\
 & React Query, Tailwind, Framer Motion, React Flow & État serveur, design, animations, graphes \\
Observabilité & Prometheus, Grafana & Métriques et tableaux de bord \\
Déploiement & Docker, Docker Compose & Conteneurisation (8 services) \\
\bottomrule
\end{tabularx}
\end{table}

\section{Technologies backend}

\subsection{FastAPI}
\textbf{Pourquoi.} Framework Python asynchrone hautes performances avec validation
native. \textbf{Comment.} Routeurs versionnés, dépendances d'autorisation,
\texttt{StreamingResponse} pour le SSE. \textbf{Bénéfices.} Concurrence par E/S
non bloquantes, documentation OpenAPI automatique. \textbf{Limites.} Écosystème
plus jeune que des frameworks synchrones établis. \textbf{Alternatives.} Django
REST Framework, Flask~\cite{fastapi}.

\subsection{PostgreSQL}
\textbf{Pourquoi.} SGBD relationnel robuste et transactionnel, idéal pour l'état et
les traces structurées. \textbf{Comment.} Onze tables, clés UUID, support JSON
pour les champs semi-structurés (\emph{previews} d'agents). \textbf{Bénéfices.}
Intégrité, requêtes analytiques, maturité. \textbf{Limites.} Inadapté au stockage
vectoriel natif, d'où le recours à Qdrant. \textbf{Alternatives.} MySQL,
SQLite~\cite{postgres}.

\subsection{SQLAlchemy et Alembic}
\textbf{Pourquoi.} ORM expressif et migrations versionnées. \textbf{Comment.}
Modèles déclaratifs asynchrones (asyncpg), migrations Alembic. \textbf{Bénéfices.}
Typage, portabilité, évolution maîtrisée du schéma. \textbf{Limites.} Courbe
d'apprentissage de l'API asynchrone. \textbf{Alternatives.} Tortoise ORM,
SQLModel.

\subsection{Redis}
\textbf{Pourquoi.} Courtier de messages à faible latence et magasin de résultats.
\textbf{Comment.} \emph{Broker} et \emph{backend} de Celery, socle de \emph{cache}.
\textbf{Bénéfices.} Rapidité, simplicité opérationnelle. \textbf{Limites.}
Persistance mémoire à dimensionner. \textbf{Alternatives.} RabbitMQ,
Amazon SQS~\cite{redis}.

\subsection{Celery}
\textbf{Pourquoi.} Standard éprouvé du traitement de tâches asynchrones en Python.
\textbf{Comment.} Tâche \texttt{process\_document\_task} (extraction, découpage,
vectorisation, indexation), relancée jusqu'à trois fois. \textbf{Bénéfices.}
Découplage du chemin critique, montée en charge des \emph{workers}.
\textbf{Limites.} Complexité opérationnelle accrue. \textbf{Alternatives.} RQ,
Dramatiq~\cite{celery}.

\subsection{Qdrant}
\textbf{Pourquoi.} Base vectorielle spécialisée et auto-hébergeable.
\textbf{Comment.} Collection indexée en HNSW, distance cosinus, vecteurs de
dimension~768. \textbf{Bénéfices.} Recherche par plus proches voisins en temps
quasi logarithmique~\cite{malkov2018}. \textbf{Limites.} Composant additionnel à
exploiter. \textbf{Alternatives.} Milvus, Weaviate, pgvector~\cite{qdrant}.

\subsection{Ollama}
\textbf{Pourquoi.} Exécution locale de modèles d'\emph{embedding}, fonctionnement
hors-ligne. \textbf{Comment.} Modèle \texttt{nomic-embed-text}, vectorisation par
lots. \textbf{Bénéfices.} Confidentialité, absence de coût d'API pour les
\emph{embeddings}. \textbf{Limites.} Dépendance aux ressources de l'hôte.
\textbf{Alternatives.} \emph{embeddings} OpenAI, Sentence-Transformers~\cite{ollama,reimers2019}.

\subsection{Gemini}
\textbf{Pourquoi.} Génération de haute qualité avec \emph{streaming} natif.
\textbf{Comment.} \texttt{gemini-2.5-flash} (repli \texttt{gemini-2.0-flash}),
température 0,2, 1000 jetons maximum, flux SSE via \texttt{streamGenerateContent}.
\textbf{Bénéfices.} Fluidité, faible latence, repli automatique.
\textbf{Limites.} Dépendance à un service distant. \textbf{Alternatives.}
OpenAI GPT, modèles locaux via Ollama~\cite{gemini}.

\section{Technologies frontend}

\subsection{Next.js}
\textbf{Pourquoi.} \emph{App Router} React 19, rendu hybride, excellente
expérience développeur. \textbf{Comment.} Pages par route, modules fonctionnels,
chargement différé des visualisations lourdes. \textbf{Bénéfices.} Découpage
automatique des \emph{bundles}, \emph{streaming} côté client. \textbf{Limites.}
Couplage fort à l'écosystème React. \textbf{Alternatives.} Remix,
SvelteKit~\cite{nextjs}.

\subsection{TypeScript}
\textbf{Pourquoi.} Sûreté de typage à grande échelle. \textbf{Comment.} Mode
strict, zéro \texttt{any} sur l'ensemble du code applicatif. \textbf{Bénéfices.}
Détection précoce des erreurs, refactorisations sûres. \textbf{Limites.} Surcoût
de typage initial. \textbf{Alternatives.} JavaScript + JSDoc, Flow.

\subsection{React Query}
\textbf{Pourquoi.} Gestion robuste de l'état serveur. \textbf{Comment.} Requêtes
mises en cache, dédupliquées et invalidées (inspecteur de traces, observatoire,
tableau de bord personnel). \textbf{Bénéfices.} Suppression des conditions de
course, standardisation des états de chargement et d'erreur. \textbf{Limites.}
Migration progressive depuis les appels manuels. \textbf{Alternatives.} SWR,
RTK Query.

\subsection{Tailwind CSS et système de design « Horizon »}
\textbf{Pourquoi.} Cohérence visuelle et productivité. \textbf{Comment.} Jetons CSS
(thème clair/sombre) mappés sur des utilitaires sémantiques ; primitives
réutilisables (Button, Card, DataCard, MetricCard, Badge, Dialog). \textbf{Bénéfices.}
Homogénéité, accessibilité, faible poids. \textbf{Limites.} Verbosité des classes.
\textbf{Alternatives.} CSS Modules, styled-components, bibliothèque shadcn/ui.

\subsection{Framer Motion et React Flow}
\textbf{Pourquoi.} Animations fluides et graphes interactifs pour les
fonctionnalités phares. \textbf{Comment.} Framer Motion anime l'AI Execution
Studio (pipeline, avatar, jauges) ; React Flow dessine l'Architecture Explorer.
\textbf{Bénéfices.} 60~images/s via animations sur \texttt{transform}/\texttt{opacity},
respect de \emph{prefers-reduced-motion}. \textbf{Limites.} Poids des
bibliothèques, atténué par le chargement différé. \textbf{Alternatives.} GSAP,
react-spring ; D3, Cytoscape.

\section{DevOps}
\textbf{Docker et Docker Compose.} L'ensemble est conteneurisé en huit services
(PostgreSQL, Redis, Qdrant, \emph{backend}, \emph{worker} Celery, \emph{frontend},
Prometheus, Grafana), garantissant la reproductibilité de l'environnement.
\textbf{Prometheus et Grafana} assurent la collecte et la visualisation des
métriques système~\cite{docker}.

\section{Détails d'implémentation}

\subsection{Implémentation backend}
Les routes délèguent aux services ; les modèles SQLAlchemy déclarent les onze
tables. Le listing~\ref{lst:model} illustre un modèle représentatif.

\begin{lstlisting}[style=code,language=Python,caption={Modèle SQLAlchemy (extrait)},label={lst:model}]
class AgentRun(UUIDPrimaryKeyMixin, Base):
    __tablename__ = "agent_runs"
    message_id: Mapped[UUID] = mapped_column(ForeignKey("messages.id"))
    agent_type: Mapped[str]
    status: Mapped[str]
    latency_ms: Mapped[int | None]
    input_preview: Mapped[dict] = mapped_column(JSONB)
    output_preview: Mapped[dict | None] = mapped_column(JSONB)
    error: Mapped[str | None]
\end{lstlisting}

\subsection{Implémentation frontend}
L'état serveur est géré par React Query. Le listing~\ref{lst:rq} illustre un
\emph{hook} d'inspection de trace orchestrant des requêtes parallèles.

\begin{lstlisting}[style=code,language=JavaScript,caption={Hook React Query (extrait)},label={lst:rq}]
export function useTraceInspector(messageId, includeContent, enabled) {
  const trace = useQuery({
    queryKey: ["trace", messageId, includeContent],
    queryFn: () => getRagTrace(messageId, includeContent),
    enabled,
  });
  // requetes secondaires en parallele, degradation gracieuse
  return { trace, /* retrieval, reranking, evaluation */ };
}
\end{lstlisting}

\subsection{Implémentation API}
L'API expose les domaines authentification, espaces, documents, chat (REST + SSE),
système, administration, RAGOps et traces, sous \texttt{/api/v1}. Les
\emph{endpoints} privilégiés (administration, RAGOps, traces) sont protégés par
des dépendances de rôle.

\subsection{Implémentation base de données}
Le schéma est défini en SQLAlchemy et son évolution gérée par Alembic (schéma
initial, évaluation des réponses, RBAC d'administration). Les fragments
référencent leur point vectoriel Qdrant via \texttt{vector\_id}.

\subsection{Implémentation RAG}
Le flux agentique instancie les sept agents et journalise chaque exécution. Le
listing~\ref{lst:wf} esquisse la boucle de génération en \emph{streaming}.

\begin{lstlisting}[style=code,language=Python,caption={Génération en streaming (extrait)},label={lst:wf}]
stream_generate = getattr(provider, "stream_generate", None)
if callable(stream_generate):
    async for token in stream_generate(prompt=prompt,
                                       system_prompt=SYSTEM_PROMPT):
        if token:
            raw_parts.append(token)
            yield ("token", {"text": token})
\end{lstlisting}

\subsection{Implémentation RAGOps}
Les services \texttt{ragops\_service} et \texttt{rag\_traces\_service} agrègent la
santé Qdrant, la file documentaire, la couverture d'\emph{embedding} et les
indicateurs de qualité, exposés via \texttt{/admin/ragops/*} et
\texttt{/admin/rag-traces/quality/summary}, puis restitués par la tour de contrôle.

\subsection{Implémentation traçabilité}
À la fin de chaque requête, le message et sa lignée (exécutions d'agents, passages
récupérés, évaluation, usage) sont persistés. L'inspecteur les recharge via des
requêtes parallèles et reconstruit les neuf sections d'analyse, en masquant les
secrets dans la vue de débogage.

\section{Conclusion}
Chaque brique a été retenue à l'issue d'une analyse de compromis explicite, puis
intégrée dans une implémentation cohérente : \emph{backend} en couches,
\emph{frontend} typé et modulaire, pipeline RAG agentique traçable, et
sous-systèmes d'exploitation et de visualisation. Le chapitre suivant valide ces
choix par les tests et les résultats.

\chapter{Validation et résultats}
\label{chap:validation}

Ce chapitre présente les interfaces réalisées, puis la stratégie de validation
(tests unitaires, d'intégration, d'API, RAG, \emph{streaming} et performance),
avant une analyse critique honnête des limites et un retour d'expérience. Les
emplacements de captures d'écran sont à renseigner manuellement à partir de
l'application en fonctionnement.

\section{Présentation des interfaces}
Pour chaque page, nous précisons l'objectif, les actions utilisateur et une
analyse fonctionnelle.

\subsection{Page d'accueil (Landing)}
\screenshotfig{Figure — Page d'accueil de DevPilot AI}{fig:cap-landing}
\textbf{Objectif.} Présenter la proposition de valeur et orienter vers la
connexion. \textbf{Actions.} Lecture des sections, navigation vers
\emph{Connexion}/\emph{Inscription}. \textbf{Analyse fonctionnelle.} Vitrine
toujours sombre, structurée en sections (problème, fonctionnalités, comparaison),
qui communique la différenciation (ancrage, traçabilité, observabilité).

\subsection{Connexion}
\screenshotfig{Figure — Interface de connexion}{fig:cap-login}
\textbf{Objectif.} Authentifier l'utilisateur. \textbf{Actions.} Saisie de
l'e-mail et du mot de passe, soumission. \textbf{Analyse fonctionnelle.} Émission
d'un jeton JWT, gestion des erreurs et redirection vers le tableau de bord ;
champs accessibles avec retours de validation.

\subsection{Inscription}
\screenshotfig{Figure — Interface d'inscription}{fig:cap-register}
\textbf{Objectif.} Créer un compte. \textbf{Actions.} Saisie du nom, de l'e-mail
et du mot de passe. \textbf{Analyse fonctionnelle.} Création d'un utilisateur de
rôle \texttt{user}, mot de passe haché, sans renvoi de données sensibles.

\subsection{Tableau de bord (personnel et administrateur)}
\screenshotfig{Figure — Tableau de bord personnel}{fig:cap-dashboard}
\textbf{Objectif.} Offrir un point d'entrée adapté au rôle. \textbf{Actions.}
Créer un espace, accéder aux documents et au chat ; pour l'administrateur,
consulter la vue d'exploitation. \textbf{Analyse fonctionnelle.} Le tableau de
bord est \emph{adapté au rôle} : les utilisateurs disposent d'un espace personnel
(création d'espaces, actions rapides), les administrateurs d'une vue d'opérations.

\subsection{Espaces de travail}
\screenshotfig{Figure — Gestion des espaces de travail}{fig:cap-workspaces}
\textbf{Objectif.} Organiser les connaissances par espace isolé. \textbf{Actions.}
Créer, ouvrir et gérer un espace et ses membres. \textbf{Analyse fonctionnelle.}
L'isolation par espace conditionne l'accès aux documents et conversations.

\subsection{Documents}
\screenshotfig{Figure — Gestion documentaire}{fig:cap-documents}
\textbf{Objectif.} Téléverser et suivre les documents. \textbf{Actions.}
Glisser-déposer, téléversement par lots, suivi de statut, relance.
\textbf{Analyse fonctionnelle.} Les statuts (\texttt{queued}, \texttt{processing},
\texttt{indexed}, \texttt{failed}) reflètent l'ingestion asynchrone ; les
documents en échec peuvent être relancés.

\subsection{Chat (streaming et citations)}
\screenshotfig{Figure — Interface de chat conversationnel}{fig:cap-chat}
\textbf{Objectif.} Poser des questions et recevoir des réponses citées.
\textbf{Actions.} Saisir une question, suivre la réponse en \emph{streaming},
consulter les sources. \textbf{Analyse fonctionnelle.} La réponse s'écrit jeton
par jeton ; les citations et sources sont reliées aux passages ; un accès à la
trace est offert aux administrateurs.

\subsection{Paramètres}
\screenshotfig{Figure — Paramètres}{fig:cap-settings}
\textbf{Objectif.} Configurer le compte et consulter l'état du système.
\textbf{Actions.} Consultation et ajustement des préférences.
\textbf{Analyse fonctionnelle.} Page personnelle accessible à tout utilisateur
authentifié.

\subsection{Tableau de bord d'administration (Operations Overview)}
\screenshotfig{Figure — Vue d'exploitation administrateur}{fig:cap-admin}
\textbf{Objectif.} Synthétiser la santé de la plateforme. \textbf{Actions.}
Consulter les métriques, surveiller les anneaux de santé et les listes de
surveillance. \textbf{Analyse fonctionnelle.} Indicateurs de santé, pipeline
documentaire, qualité RAG, latence des agents, usage et coût des modèles, et
listes d'actions/réponses/documents à risque.

\subsection{Tableau de bord Qualité}
\screenshotfig{Figure — Tableau de bord Qualité}{fig:cap-quality}
\textbf{Objectif.} Analyser la qualité des réponses. \textbf{Actions.} Lire les
distributions et les listes des pires réponses. \textbf{Analyse fonctionnelle.}
Carte de score, distributions de fidélité et d'hallucination, issues des réponses
et latence par agent.

\subsection{Tour de contrôle RAGOps}
\screenshotfig{Figure — Tour de contrôle RAGOps}{fig:cap-ragops}
\textbf{Objectif.} Superviser l'exploitation du pipeline. \textbf{Actions.}
Examiner la santé des services, la santé par espace et le centre des échecs.
\textbf{Analyse fonctionnelle.} Santé Qdrant/Redis/Celery, couverture
d'\emph{embedding}, file d'ingestion et relance des documents en échec.

\subsection{Observatoire des traces RAG}
\screenshotfig{Figure — Observatoire des traces (liste)}{fig:cap-traces}
\textbf{Objectif.} Explorer et filtrer les traces. \textbf{Actions.} Filtrer
(espace, utilisateur, fidélité, hallucination, dates), trier, ouvrir une trace.
\textbf{Analyse fonctionnelle.} Cartes de trace interactives avec score de santé,
distributions et bandeau « à investiguer ».

\subsection{Détail d'une trace (AI Trace Inspector)}
\screenshotfig{Figure — Inspecteur de traces}{fig:cap-trace-detail}
\textbf{Objectif.} Reconstruire le cycle de vie complet d'une requête.
\textbf{Actions.} Parcourir les neuf sections d'analyse. \textbf{Analyse
fonctionnelle.} Résumé, chronologie des agents, explorateurs de récupération, de
reclassement et de citations, centre d'évaluation, analyse du correcteur (\emph{diff}),
usage des modèles et graphe interactif.

\subsection{AI Execution Studio}
\screenshotfig{Figure — AI Execution Studio}{fig:cap-studio}
\textbf{Objectif.} Visualiser l'exécution du pipeline en temps réel.
\textbf{Actions.} Lancer une question, observer les dix étapes animées, rejouer.
\textbf{Analyse fonctionnelle.} Pipeline animé piloté par le flux SSE réel, avatar
réactif et contrôles de lecture (1×/2×/5×).

\subsection{Architecture Explorer}
\screenshotfig{Figure — Architecture Explorer}{fig:cap-archi}
\textbf{Objectif.} Explorer l'architecture de façon interactive. \textbf{Actions.}
Lancer une requête (flux animé), cliquer un n{\oe}ud pour l'inspecter.
\textbf{Analyse fonctionnelle.} Carte React Flow avec flux de requête animé,
inspecteur de n{\oe}ud (rôle, configuration réelle, métriques d'exécution) et
pages immersives par composant.

\section{Tests unitaires}
La couche de présentation est couverte par \textbf{39 tests} (Vitest et Testing
Library) portant sur les fonctions pures : algorithme de \emph{diff} mot à mot
(plus longue sous-séquence commune), score de santé de trace, mappage des tons,
construction des paramètres de filtre de l'observatoire, mise en \emph{buckets}
des distributions, et rendu des composants \texttt{StatusBadge}/\texttt{Badge}. Le
\emph{backend} dispose d'une suite \texttt{pytest}.

\screenshotfig{Figure — Exécution de la suite de tests unitaires}{fig:cap-tests}

\section{Tests d'intégration}
Les tests d'intégration valident la chaîne ingestion $\rightarrow$ indexation
$\rightarrow$ récupération sur un corpus contrôlé, en vérifiant que les fragments
téléversés deviennent récupérables avec des scores cohérents.

\section{Tests d'API}
Les contrats des \emph{endpoints} sont vérifiés (codes de statut, schémas,
autorisation) à l'aide de requêtes \texttt{curl} contre l'API réelle, en
particulier le rejet des accès non autorisés (401/403) et la non-divulgation des
secrets par \texttt{/system/ai-config}.

\section{Tests RAG}
La qualité est validée selon trois axes : récupération (pertinence des passages),
ancrage (réponse étayée par le contexte) et citation (références correctes).
L'agent évaluateur fournit une mesure continue de la fidélité, de la pertinence,
de la précision du contexte et du risque d'hallucination.

\section{Tests de streaming}
La diffusion SSE a été vérifiée par \texttt{curl} : les événements arrivent
\emph{progressivement} (\texttt{trace}, puis une succession de \texttt{token},
puis \texttt{evaluation}, \texttt{message}, \texttt{done}), confirmant un
\emph{streaming} réel de jetons après la correction du comportement de génération
simulée du fournisseur.

\section{Tests de performance}
Les routes pèsent de 110 à 175~kB environ ; les visualisations lourdes sont
chargées de façon différée ; les animations s'exécutent à 60~images/s grâce à
l'usage exclusif de \texttt{transform}/\texttt{opacity} ; le \emph{backend}
asynchrone et la recherche HNSW concourent à une expérience réactive. Le
tableau~\ref{tab:accept} synthétise les critères d'acceptation.

\begin{table}[H]
\centering\caption{Critères d'acceptation}\label{tab:accept}\small
\begin{tabularx}{\textwidth}{@{}X X c@{}}
\toprule
\textbf{Objet validé} & \textbf{Critère} & \textbf{Résultat} \\
\midrule
Compilation \emph{frontend} & \texttt{next build} sans erreur & \checkmark \\
Qualité du code & \texttt{eslint} sans avertissement & \checkmark \\
Tests unitaires & 39 tests au vert & \checkmark \\
Streaming & Jetons progressifs (SSE) & \checkmark \\
Réponse RAG & Ancrée et citée & \checkmark \\
Évaluation & Calculée et persistée & \checkmark \\
Traçabilité & Trace reconstructible & \checkmark \\
Sécurité & Aucune fuite de secret / inter-espaces & \checkmark \\
Tests e2e automatisés (Playwright) & Couverture des parcours & À venir \\
\bottomrule
\end{tabularx}
\end{table}

\section{Analyse critique}
Le système atteint un niveau de maturité notable, mais plusieurs points restent
perfectibles. La migration de la couche de données vers React Query est en cours
sur certaines pages historiques. L'évaluation repose sur des heuristiques et un
juge LLM plutôt que sur un cadre formel à large échelle. Le déploiement cible est
mono-machine. Enfin, la couverture de tests automatisés gagnerait à être étendue
aux parcours de bout en bout.

\section{Limitations du projet}
La comptabilité fine d'usage des modèles dépend des métadonnées exposées par
chaque fournisseur et n'est pas toujours disponible ; l'inspecteur affiche alors
honnêtement « non capturé ». Le multi-tenant reste au niveau de l'espace de
travail, sans cloisonnement organisationnel complet. Les visualisations de
« l'univers vectoriel » représentent des positions projetées et non les
coordonnées réelles en 768 dimensions, choix de visualisation assumé.

\section{Retour d'expérience}
Ce projet a confirmé que la valeur d'une plateforme d'IA d'entreprise ne réside
pas seulement dans la qualité des réponses, mais dans leur auditabilité et leur
exploitabilité. Concevoir pour l'observabilité dès l'origine, traiter les
exigences non fonctionnelles comme des exigences de premier ordre et raisonner
systématiquement en compromis ont été les enseignements majeurs, aux côtés de la
maîtrise opérationnelle d'une pile distribuée conteneurisée.

\section{Conclusion}
La stratégie de validation couvre l'ensemble des couches, du test unitaire à
l'acceptation, en accordant une attention particulière aux propriétés spécifiques
d'une plateforme d'IA : ancrage, citations, qualité mesurée, \emph{streaming} et
sécurité. Les résultats confirment que \devpilot{} satisfait ses objectifs, tout
en identifiant honnêtement ses axes d'amélioration.

%====================== CONCLUSION GÉNÉRALE ======================
\chapter*{Conclusion générale}
\addcontentsline{toc}{chapter}{Conclusion générale}
\markboth{Conclusion générale}{Conclusion générale}

Ce projet de fin d'études a abouti à la conception et à la réalisation de
\devpilot{}, une plateforme agentique de RAG de qualité production. L'objectif
initial — combler l'écart entre la puissance générative des grands modèles de
langage et l'exigence de réponses \emph{fondées}, \emph{citées}, \emph{mesurées}
et \emph{traçables} — est atteint et validé sur l'ensemble du pipeline.

\paragraph{Réalisations.} La plateforme centralise les connaissances en espaces
isolés, les ingère de façon asynchrone, les indexe vectoriellement, et restitue
des réponses ancrées et citées via un flux agentique de sept agents diffusé en
\emph{streaming}. Chaque réponse est évaluée et, le cas échéant, corrigée.
L'ensemble est instrumenté par trois contributions distinctives : la tour de
contrôle \textbf{RAGOps}, le sous-système de \textbf{traçabilité} (explorateur et
inspecteur de traces) et l'\textbf{AI Execution Studio}.

\paragraph{Maturité architecturale.} L'architecture en couches, asynchrone et
neutre vis-à-vis des fournisseurs d'IA, la discipline de typage statique de bout
en bout, la conteneurisation complète et la fondation de tests confèrent au
système une maturité dépassant celle d'un prototype académique.

\paragraph{Compétences acquises.} Ce travail a consolidé des compétences en
architecture des systèmes distribués, en ingénierie des plateformes d'IA (RAG,
agents, évaluation), en observabilité et traçabilité, ainsi qu'en ingénierie
\emph{frontend} avancée et en pratiques DevOps.

\paragraph{Enseignement principal.} La valeur d'une plateforme d'IA d'entreprise
ne réside pas seulement dans la qualité de ses réponses, mais dans leur
\emph{auditabilité} et leur \emph{exploitabilité}. Concevoir pour l'observabilité
dès l'origine et raisonner systématiquement en compromis d'ingénierie ont été les
leviers déterminants de ce projet.

%====================== TRAVAUX FUTURS ======================
\chapter*{Perspectives et travaux futurs}
\addcontentsline{toc}{chapter}{Perspectives et travaux futurs}
\markboth{Perspectives}{Perspectives}

Plusieurs axes d'évolution prolongent naturellement ce travail.

\begin{description}[leftmargin=2.4cm,style=nextline]
\item[Orchestration multi-agents] Adoption d'un graphe d'exécution explicite
(type LangGraph) pour modéliser les transitions, les boucles et les stratégies de
repli du flux agentique.
\item[RAG avancé] Récupération hybride enrichie, fenêtre de contexte adaptative,
et \textbf{GraphRAG} exploitant un graphe de connaissances pour les questions
multi-sauts.
\item[Évaluation formelle] Intégration de RAGAS~\cite{es2023} et de jeux
d'évaluation versionnés pour un suivi longitudinal de la qualité.
\item[MCP et connecteurs d'entreprise] Intégration du \emph{Model Context
Protocol} et de connecteurs (dépôts de code, wikis, tickets) pour étendre les
sources de connaissance.
\item[Agents autonomes] Agents capables de planifier des actions multi-étapes et
d'invoquer des outils sous contrôle de gouvernance.
\item[RAGOps avancé et gouvernance de l'IA] Détection d'anomalies, budgets de
coût et de latence, politiques de rétention, et pistes d'audit étendues pour la
gouvernance de l'IA.
\item[Déploiement cloud et multi-tenant] Migration vers une orchestration
Kubernetes, cloisonnement organisationnel, quotas et facturation par locataire.
\end{description}

Ces perspectives s'appuient sur les fondations posées — traçabilité native,
architecture découplée et observabilité — et visent à faire évoluer \devpilot{}
d'une plateforme de démonstration aboutie vers un produit d'entreprise à grande
échelle.


%------------------------------------------------------------------ Bibliographie
\begin{thebibliography}{99}
\addcontentsline{toc}{chapter}{Bibliographie}

\bibitem{vaswani2017} A.~Vaswani \emph{et al.}, « Attention Is All You Need », \emph{NeurIPS}, 2017.
\bibitem{lewis2020} P.~Lewis \emph{et al.}, « Retrieval-Augmented Generation for Knowledge-Intensive NLP Tasks », \emph{NeurIPS}, 2020.
\bibitem{karpukhin2020} V.~Karpukhin \emph{et al.}, « Dense Passage Retrieval for Open-Domain Question Answering », \emph{EMNLP}, 2020.
\bibitem{robertson2009} S.~Robertson, H.~Zaragoza, « The Probabilistic Relevance Framework: BM25 and Beyond », \emph{Found. Trends IR}, 2009.
\bibitem{malkov2018} Y.~Malkov, D.~Yashunin, « Efficient and Robust ANN Search Using HNSW Graphs », \emph{IEEE TPAMI}, 2018.
\bibitem{reimers2019} N.~Reimers, I.~Gurevych, « Sentence-BERT », \emph{EMNLP}, 2019.
\bibitem{nogueira2019} R.~Nogueira, K.~Cho, « Passage Re-ranking with BERT », \emph{arXiv:1901.04085}, 2019.
\bibitem{es2023} S.~Es \emph{et al.}, « RAGAS: Automated Evaluation of Retrieval Augmented Generation », \emph{arXiv:2309.15217}, 2023.
\bibitem{fastapi} S.~Ramírez, \emph{FastAPI Documentation}, \url{https://fastapi.tiangolo.com}, 2026.
\bibitem{postgres} \emph{PostgreSQL 16 Documentation}, \url{https://www.postgresql.org/docs/16}, 2026.
\bibitem{qdrant} \emph{Qdrant Vector Database Documentation}, \url{https://qdrant.tech/documentation}, 2026.
\bibitem{celery} \emph{Celery — Distributed Task Queue}, \url{https://docs.celeryq.dev}, 2026.
\bibitem{redis} \emph{Redis Documentation}, \url{https://redis.io/docs}, 2026.
\bibitem{nextjs} \emph{Next.js Documentation}, \url{https://nextjs.org/docs}, 2026.
\bibitem{docker} \emph{Docker Documentation}, \url{https://docs.docker.com}, 2026.
\bibitem{gemini} Google DeepMind, \emph{Gemini API Documentation}, \url{https://ai.google.dev}, 2026.
\bibitem{ollama} \emph{Ollama Documentation}, \url{https://ollama.com}, 2026.
\bibitem{langsmith} LangChain, \emph{LangSmith — LLM Application Observability}, \url{https://docs.smith.langchain.com}, 2026.
\bibitem{llamaindex} \emph{LlamaIndex Documentation}, \url{https://docs.llamaindex.ai}, 2026.
\bibitem{haystack} deepset, \emph{Haystack Documentation}, \url{https://haystack.deepset.ai}, 2026.

\end{thebibliography}

\end{document}
